% LABYRINTHOGRAPHY — a zine. A5, 28 pages, 1-bit black.
% Build: make zine   (needs `make` first: every number and figure comes from build/)
\documentclass[9pt,twoside]{memoir}
\setstocksize{210mm}{148mm}
\settrimmedsize{\stockheight}{\stockwidth}{*}
\setlrmarginsandblock{14mm}{14mm}{*}
\setulmarginsandblock{14mm}{22mm}{*}
\setheadfoot{0mm}{10mm}
\checkandfixthelayout
\usepackage{fontspec}
\defaultfontfeatures{Ligatures=TeX}
\setmainfont{EBGaramond}[Path=../cache/fonts/, Extension=.otf, UprightFont=*-Regular,
  ItalicFont=*-Italic, BoldFont=*-SemiBold, Numbers=OldStyle]
\newfontfamily\figfont{EBGaramond}[Path=../cache/fonts/, Extension=.otf, UprightFont=*-Regular,
  ItalicFont=*-Italic, BoldFont=*-SemiBold, Numbers={Lining,Monospaced}]
\newfontfamily\monofont{IBMPlexMono}[Path=../cache/fonts/, Extension=.ttf, UprightFont=*-Regular,
  BoldFont=*-Medium, LetterSpace=6]
\usepackage{microtype}
\usepackage{graphicx}
\usepackage{tikz}
\usepackage{booktabs}
\usepackage{amsmath}
\setlength\parindent{0pt}
\setlength\parskip{0.6em}
\linespread{1.08}
\raggedright
\pagestyle{empty}
% numbers from the engine
\newcommand\lab[1]{\csname lab@#1\endcsname}
\input{../build/the-trial/numbers.tex}
\input{../build/compare/numbers.tex}
\input{../build/zine/first-steps.tex}
% --- the page grammar: a kicker at the top, the image in the middle, a few lines below
\newcommand\mono[1]{{\monofont\fontsize{6.4}{8}\selectfont\MakeUppercase{#1}}}
\newcommand\folio{\vfill\noindent\makebox[\textwidth][\ifodd\value{page}r\else l\fi]{\mono{\thepage}}}
\newcommand\kicker[2]{\noindent\mono{#1\quad —\quad #2}\par\vspace{6mm}}
\newcommand\zpage[3]{% #1 number, #2 title, #3 body (figure and text)
  \kicker{#1}{#2}#3\folio\newpage}
\newcommand\zcap[1]{\par\vspace{5mm}{\fontsize{9.4}{12.4}\selectfont #1}\par}
\newcommand\zsmall[1]{\par\vspace{2mm}{\monofont\fontsize{5.6}{7}\selectfont #1}\par}
\newcommand\bk[1]{\emph{\lab{c-#1-title}}}
\newcommand\cdir{../build/compare/}
% the maze labels and record macros used by generated files
\newcommand\lblhead[1]{\MakeUppercase{#1}}
\newcommand\lblletter[1]{#1}
\tikzset{mazelabel/.style={inner sep=0pt, font=\monofont\fontsize{5.4}{6.6}\selectfont},
         shaftlabel/.style={inner sep=0.5pt, fill=white, anchor=west, font=\monofont\fontsize{5}{6}\selectfont}}
\newcommand\labelledfigure[2]{\begin{tikzpicture}[x=1mm,y=1mm]
  \node[anchor=north west, inner sep=0pt] at (0,0) {\includegraphics{#1}}; \input{#2}\end{tikzpicture}}
\newcommand\vocabdisc[2]{\begin{tikzpicture}
  \node[inner sep=0pt] (vd) {\includegraphics[width=#2]{#1}};
  \draw[line width=0.2pt] (vd.center) circle ({\lab{vocab-rim}*#2});\end{tikzpicture}}
\newcommand\five[2]{% #1 prefix, #2 width
  \begin{center}\begin{tabular}{@{}*{5}{c}@{}}
  \includegraphics[width=#2]{\cdir #1-the-trial.pdf} & \includegraphics[width=#2]{\cdir #1-bartleby.pdf} &
  \includegraphics[width=#2]{\cdir #1-yellow-wallpaper.pdf} & \includegraphics[width=#2]{\cdir #1-heart-of-darkness.pdf} &
  \includegraphics[width=#2]{\cdir #1-pride-and-prejudice.pdf}\\[1mm]
  \mono{Trial} & \mono{Bartleby} & \mono{Wallpaper} & \mono{Darkness} & \mono{Pride}
  \end{tabular}\end{center}}
\newcommand\strip[2]{\noindent\mono{#2}\par\vspace{0.8mm}\includegraphics[width=\textwidth,height=16.5mm,keepaspectratio]{\cdir #1}\par\vspace{1.8mm}}

\begin{document}

% ------------------------------------------------------------------ 1 cover
\noindent\mono{Labyrinthography}\hfill\mono{01}
\vfill
\begin{center}\includegraphics{../build/zine/cover-maze.pdf}\end{center}
\vfill
\noindent\mono{Novels drawn as the walk their sentences make}
\newpage

% ------------------------------------------------------------------ 2
\zpage{01}{The walk}{%
{\fontsize{12}{16}\selectfont Read a novel one sentence at a time, and walk.\par}
\vspace{4mm}
Every sentence is a straight line, as long as it has words. Before each line,
turn. A warm sentence turns you a quarter to the left. A cold one turns you to
the right. A plain one, a door, a staircase, keeps you going straight.

How a sentence feels is counted with a published list of some 3,400 English
words, each scored from $-5$ to $+5$. Names are not feelings, so \emph{Miss}
and \emph{Block} are not counted.

Nothing is placed by hand. The drawing is the book.
}

% ------------------------------------------------------------------ 3
\zpage{02}{Eight sentences}{%
\vfill
\begin{center}\includegraphics{../build/zine/first-steps.pdf}\end{center}
\vfill
\zcap{K. waits up for Miss Bürstner (\stepsfrom--\stepsto). ``Oh, Mr.~K.!''
turns left. ``It's me,'' said K., and the path goes straight. ``It is a little
odd, isn't it?'' turns right.}
}

% ------------------------------------------------------------------ 4
\zpage{03}{The Trial, Franz Kafka}{%
The whole novel, \lab{sentences} sentences, opposite. It begins at the black
square and ends at the open one.

\vspace{4mm}
{\monofont\fontsize{6.4}{10}\selectfont
A \quad KNOT I\quad the densest tangle\\
B \quad KNOT II\\
C \quad HUB\quad the most-crossed sentence\\
D \quad RETURN\quad two sentences \lab{ret-gap} apart, crossing\\
E \quad COIL\quad the tightest spiral\\
F \quad DARKEST\quad thirty cold sentences in a row\\
G \quad CORRIDOR\quad the longest plain stretch\\
H \quad PALIMPSEST\quad a sentence drawn over another\par}

\vspace{4mm}
The tangles sit in the middle of the book, where K. goes from the lawyer to the
painter to Block. The places he goes to learn about his case, hundreds of pages
apart, end up on the same small patch of ground.
}

% ------------------------------------------------------------------ 5 the maze
\vspace*{-4mm}
\noindent\hspace*{-1mm}\scalebox{0.93}{\labelledfigure{../build/the-trial/maze.pdf}{../build/the-trial/maze-labels.tex}}
\folio\newpage

% ------------------------------------------------------------------ 6
\zpage{04}{C — The hub}{%
\vfill
\begin{center}\includegraphics[width=80mm]{../build/the-trial/cut-C.pdf}\end{center}
\vfill
\zcap{The sentence the path crosses most, \lab{hub-n} times. A court official:
``\lab{hub-quote}'' Most of the lines that cut across it are Block, describing
years of hearings that ``went through the same things anyway''.}
}

% ------------------------------------------------------------------ 7
\zpage{05}{E — The coil}{%
\vfill
\begin{center}\includegraphics[width=80mm]{../build/the-trial/cut-E.pdf}\end{center}
\vfill
\zcap{Two and a half full circles in twelve sentences: the lawyer reporting
progress. ``The results are very encouraging.'' The word list hears hope. The
novel knows better: ``there was always some progress, but he could never be
told what sort of progress it was.''}
}

% ------------------------------------------------------------------ 8
\zpage{06}{Feelings that linger}{%
{\fontsize{12}{16}\selectfont After a cold sentence, Kafka writes another cold one
\lab{pp-R}\,\% of the time. Chance would give \lab{base-R}\,\%.\par}
\vspace{6mm}
Shuffle the novel's sentences like cards, two hundred times, and draw each
shuffle. Every shuffle has the same sentences. Only the order is lost.

Not one of the two hundred kept its moods together the way the novel does.
Feeling in \emph{The Trial} does not strike. It settles over a passage and
stays, and the path folds back on itself.

Everything else (how big the drawing is, how far it wanders, how often it
crosses) the shuffles do too. Kafka's order decides only where the tangle sits.
}

% ------------------------------------------------------------------ 9
\zpage{07}{The shaft, and one word}{%
\begin{center}
\includegraphics[width=50mm]{../build/the-trial/mineshaft-half.pdf}\hspace{6mm}%
\includegraphics[width=50mm]{../build/the-trial/mineshaft-r.pdf}\\[1mm]
\makebox[50mm]{\mono{every word}}\hspace{6mm}\makebox[50mm]{\mono{without "like"}}
\end{center}
\zcap{Build the walk in three dimensions: warm sentences tilt it up, cold
ones down. Scored as the list says, \emph{The Trial} climbs. But its largest
source of warmth is \emph{like}, used mostly as a comparison (``like a dog''),
not as affection. Leave that one word out and the book sinks. A drawing can show
the limits of its own method.}
}

% ------------------------------------------------------------------ 10
\zpage{08}{Standing among the sentences}{%
\vfill
\includegraphics[width=\textwidth]{../build/the-trial/variants/quarter-high.png}
\vfill
\zcap{The busiest junction of the walk, seen from above. Every sentence stands
on its own line as flat black letters. Where the walk crosses itself, sentences
pass through each other.}
}

% ------------------------------------------------------------------ 11
\zpage{09}{Walked word by word}{%
\vfill
\includegraphics[width=\textwidth]{../build/the-trial/variants/rolling-high.png}
\vfill
\zcap{The same walk, turning at every word by the mood of the last twelve. The
corners melt, and each sentence stands along a curve.}
}

% ------------------------------------------------------------------ 12
\zpage{10}{Five books}{%
{\fontsize{12}{16}\selectfont A drawing says little until there is another beside it.\par}
\vspace{5mm}
\textsc{The Trial}, Franz Kafka, \lab{c-the-trial-year}\\
\textsc{Bartleby, the Scrivener}, Herman Melville, \lab{c-bartleby-year}\\
\textsc{The Yellow Wallpaper}, Charlotte Perkins Gilman, \lab{c-yellow-wallpaper-year}\\
\textsc{Heart of Darkness}, Joseph Conrad, \lab{c-heart-of-darkness-year}\\
\textsc{Pride and Prejudice}, Jane Austen, \lab{c-pride-and-prejudice-year}

\vspace{5mm}
An office and a refusal. A room and a journal. A river and a voice. A drawing
room and a marriage. Each answers \emph{The Trial} in some way, and each is
scored, walked and shuffled exactly the same.
}

% ------------------------------------------------------------------ 13
\zpage{11}{Five walks}{%
\vfill
\five{shape}{21.5mm}
\vspace{6mm}
\five{scale}{21.5mm}
\vfill
\zcap{Above, each in its own frame: only the shape. Below, all at one scale:
the size. \bk{yellow-wallpaper} is a pocket of \lab{c-yellow-wallpaper-sentences}
sentences; \bk{pride-and-prejudice} runs to \lab{c-pride-and-prejudice-sentences}.
Austen leans warm by a wide margin; Conrad and Melville lean cold.}
}

% ------------------------------------------------------------------ 14
\zpage{12}{Vocabulary records}{%
\begin{center}
\vocabdisc{\cdir vocab-the-trial.png}{34mm}\hfill\vocabdisc{\cdir vocab-bartleby.png}{34mm}\hfill\vocabdisc{\cdir vocab-yellow-wallpaper.png}{34mm}\\
\makebox[34mm]{\mono{Trial}}\hfill\makebox[34mm]{\mono{Bartleby}}\hfill\makebox[34mm]{\mono{Wallpaper}}\\[3mm]
\vocabdisc{\cdir vocab-heart-of-darkness.png}{34mm}\hspace{8mm}\vocabdisc{\cdir vocab-pride-and-prejudice.png}{34mm}\\
\makebox[34mm]{\mono{Darkness}}\hspace{8mm}\makebox[34mm]{\mono{Pride}}
\end{center}
\zcap{The dictionary as the groove of a record, \lab{vocab-words} words: the
commonest at the centre, the rarest at the rim. Every word a book uses is a
black dot. Stretch for stretch, \emph{The Trial} has the plainest vocabulary of the
five, \lab{c-the-trial-types-6k} words in 6,000; Conrad the richest,
\lab{c-heart-of-darkness-types-6k}.}
}

% ------------------------------------------------------------------ 15
\zpage{13}{Tree rings}{%
\vfill
\five{rings}{21.5mm}
\vfill
\zcap{Cut a book like a trunk. The first page is the centre, each chapter a
ring as large as it is long, each sentence a grain as long as its words.
\emph{Pride and Prejudice} grows in \lab{c-pride-and-prejudice-rings} thin rings;
\emph{Heart of Darkness} in three broad seasons.}
}

% ------------------------------------------------------------------ 16
\zpage{14}{The loom}{%
\five{loom}{21.5mm}
\zcap{The names a book uses most are warp threads; the book runs down
through them. K. is a solid black thread, there on every stretch. Elizabeth is in
\lab{c-pride-and-prejudice-lead-share}\,\% of hers. Kurtz, the man \emph{Heart of
Darkness} is about, is in only \lab{c-heart-of-darkness-lead-share}\,\%.}
}

% ------------------------------------------------------------------ 17
\zpage{15}{Echoes}{%
\strip{echo-the-trial.pdf}{Trial}
\strip{echo-bartleby.pdf}{Bartleby}
\strip{echo-yellow-wallpaper.pdf}{Wallpaper}
\strip{echo-heart-of-darkness.pdf}{Darkness}
\strip{echo-pride-and-prejudice.pdf}{Pride}
\zcap{Every recurring word is one arch, spanning the longest time the book
forgets it. In \emph{The Yellow Wallpaper} the longest are \emph{yellow} and
\emph{moonlight}.}
}

% ------------------------------------------------------------------ 18
\zpage{16}{Tides of speech}{%
\strip{tide-the-trial.pdf}{Trial}
\strip{tide-bartleby.pdf}{Bartleby}
\strip{tide-yellow-wallpaper.pdf}{Wallpaper}
\strip{tide-heart-of-darkness.pdf}{Darkness}
\strip{tide-pride-and-prejudice.pdf}{Pride}
\zcap{The share of words spoken inside quotation marks. \emph{Heart of
Darkness} is a flood, \lab{c-heart-of-darkness-speech}\,\%: the whole tale is
Marlow talking. \emph{The Yellow Wallpaper}, a secret journal, is nearly dry.}
}

% ------------------------------------------------------------------ 19
\zpage{17}{Breath}{%
\strip{breath-the-trial.pdf}{Trial}
\strip{breath-bartleby.pdf}{Bartleby}
\strip{breath-yellow-wallpaper.pdf}{Wallpaper}
\strip{breath-heart-of-darkness.pdf}{Darkness}
\strip{breath-pride-and-prejudice.pdf}{Pride}
\zcap{Every clause, from one comma to the next, as tall as it is long.
\emph{The Trial} holds its breath longest: \lab{c-the-trial-longbreath}\,\% of its
clauses run to twenty words or more.}
}

% ------------------------------------------------------------------ 20
\zpage{18}{The punctuation sky}{%
\newcommand\skyrow[2]{\noindent\mono{#2}\par\vspace{0.8mm}\includegraphics[width=\textwidth]{\cdir #1}\par\vspace{2.2mm}}
\skyrow{sky-the-trial.pdf}{Trial}
\skyrow{sky-bartleby.pdf}{Bartleby}
\skyrow{sky-yellow-wallpaper.pdf}{Wallpaper}
\skyrow{sky-heart-of-darkness.pdf}{Darkness}
\zcap{One star per sentence: a dot for a full stop, a ring for a question, a bar
for an exclamation. \emph{The Yellow Wallpaper} cries out in
\lab{c-yellow-wallpaper-exclaim}\,\% of its sentences. Ten times \emph{The Trial}.}
}

% ------------------------------------------------------------------ 21
\zpage{19}{The pronoun walk}{%
\vfill
\five{pronouns}{21.5mm}
\vfill
\zcap{Walk again, steered by person: \emph{I} north, \emph{you} south,
\emph{he} east, \emph{she} west. \emph{The Trial} walks due east: it is told
about K., never by him. \emph{The Yellow Wallpaper} walks due north. Only
\emph{Pride and Prejudice} turns west.}
}

% ------------------------------------------------------------------ 22
\zpage{20}{The clock}{%
\vfill
\five{clock}{21.5mm}
\vfill
\zcap{Every time of day a book mentions, on a 24-hour dial, midnight at the top;
further out, later in the book. \emph{Bartleby} keeps office hours, around
\lab{c-bartleby-hour}. \emph{The Yellow Wallpaper} lives at night, around
\lab{c-yellow-wallpaper-hour}.}
}

% ------------------------------------------------------------------ 23
\zpage{21}{The light meter}{%
\noindent\mono{Trial}\par\vspace{1.2mm}\includegraphics[width=\textwidth]{\cdir light-the-trial.png}\par\vspace{3mm}
\noindent\mono{Bartleby}\par\vspace{1.2mm}\includegraphics[width=\textwidth]{\cdir light-bartleby.png}\par\vspace{3mm}
\noindent\mono{Wallpaper}\par\vspace{1.2mm}\includegraphics[width=\textwidth]{\cdir light-yellow-wallpaper.png}\par\vspace{3mm}
\noindent\mono{Darkness}\par\vspace{1.2mm}\includegraphics[width=\textwidth]{\cdir light-heart-of-darkness.png}\par\vspace{3mm}
\noindent\mono{Pride}\par\vspace{1.2mm}\includegraphics[width=\textwidth]{\cdir light-pride-and-prejudice.png}\par
\zcap{Light words (\emph{sun, lamp, candle, white}) against dark ones
(\emph{shadow, fog, night, black}), read along each book. \emph{Heart of Darkness}
is the most marked. \emph{The Trial} is lighter than its name: candles carried
through dim rooms. Austen hardly mentions light at all.}
}

% ------------------------------------------------------------------ 24
\zpage{22}{Other ways to walk}{%
\vfill
\begin{center}
\includegraphics[width=52mm]{../build/the-trial/variants/graded-plan.pdf}\hspace{6mm}%
\includegraphics[width=52mm]{../build/the-trial/variants/rolling-plan.pdf}\\[1mm]
\makebox[52mm]{\mono{graded}}\hspace{6mm}\makebox[52mm]{\mono{rolling}}
\end{center}
\vfill
\zcap{\emph{The Trial} again. Graded: each sentence turns as far as it feels,
and the maze becomes a coastline. Rolling: the walker turns at every word, by
the mood of the last twelve, and the corners melt into curves.}
}

% ------------------------------------------------------------------ 25
\zpage{23}{Five portraits}{%
\textsc{The Trial} walks evenly, in quiet moods that linger, with the plainest
words, around one man who is never absent. Its tangles sit in the rooms where he
looks for answers.

\textsc{Bartleby} is short and cool, and asks the most questions. It keeps office
hours.

\textsc{The Yellow Wallpaper} is the smallest walk: few names, almost no other
voices, written at night, and exclaiming.

\textsc{Heart of Darkness} is one long spoken voice, leaning cold, rich in words,
built around a man who is mostly not there.

\textsc{Pride and Prejudice} is the largest walk and the warmest: half spoken,
densely woven around Elizabeth, and the only one that walks west.
}

% ------------------------------------------------------------------ 26
\zpage{24}{How this was made}{%
\fontsize{8.6}{11.4}\selectfont
Texts from Project Gutenberg. Sentences split at full stops, question and
exclamation marks, never after an initial (\emph{Josef K.}).
Feeling: the AFINN-165 word list (F.\,Å. Nielsen, 2011), summed per sentence.
Capitalised words in mid-sentence are treated as names and not scored.

Every claim about order is tested against two hundred shuffles of the same
sentences, and only called a property of a book when it stands clearly apart
from all of them. Every test is also repeated without the ambiguous word
\emph{like}.

A word list does not hear irony, does not notice \emph{not}, and for a translated
novel it reads the translator. The drawings show the surface of the words, in
the order the author chose. That is all they claim.

Rarity of words: the system dictionary, ordered by the \texttt{wordfreq} corpus.
Every number in this zine is computed by the program that draws it.
}

% ------------------------------------------------------------------ 27 colophon
\kicker{25}{Colophon}
\vfill
{\monofont\fontsize{6}{9}\selectfont
LABYRINTHOGRAPHY\\
A ZINE\\[3mm]
SET IN EB GARAMOND AND IBM PLEX MONO\\
PRINTED IN BLACK ONLY\\
A5, SADDLE-STITCHED\\[3mm]
CODE: GITHUB.COM/SLOEV/TYPOGRAPHY\_MAZE\\[3mm]
THE TRIAL QUOTED IN THE TRANSLATION BY DAVID WYLLIE,\\
PROJECT GUTENBERG EBOOK 7849 (COPYRIGHTED; QUOTED AND\\
DRAWN BY PERMISSION ONLY). OTHER TEXTS PUBLIC DOMAIN.\par}
\folio\newpage

% ------------------------------------------------------------------ 28 back cover
\vspace*{\fill}
\begin{center}\vocabdisc{\cdir vocab-the-trial.png}{46mm}\end{center}
\vspace{6mm}
\begin{center}\mono{Every dot is a word The Trial uses}\end{center}
\vspace*{\fill}
\noindent\mono{Labyrinthography}\hfill\mono{01}

\end{document}
